\pdfvariable suppressoptionalinfo 611
\providecommand{\CRevisionRound}{2}
\documentclass[10pt]{article}
\usepackage[a4paper,margin=18mm]{geometry}
\usepackage{amsmath,amssymb,amsthm,booktabs,microtype}
\usepackage[T1]{fontenc}
\usepackage{lmodern}
\usepackage[hidelinks,hypertexnames=false]{hyperref}
\newtheorem{theorem}{Theorem}
\newtheorem{proposition}{Proposition}
\newtheorem{lemma}{Lemma}
\title{Exact Itinerary Intervals and Mode Locking for a Contracted Rotation}
\author{HCS Research Program}
\date{30 August 2026\quad (revision \CRevisionRound)}

\begin{document}
\maketitle
\begin{abstract}
We give a finite, exact atlas for the discontinuous contracted rotation
\(f_{\lambda,\delta}(x)=\{\lambda x+\delta\}\) on \([0,1)\).  A binary carry
word determines an affine return map with slope \(\lambda^n\), hence one
candidate fixed point.  For \(\lambda\in\{1/2,2/3,3/4\}\) we enumerate all
747 primitive cyclic representatives through length 12 for each slope and
intersect the half-open branch inequalities using rational arithmetic.  The
receipt contains 2241 word rows, 138 nonempty components, exact endpoint
equalities, and 295 independent direct-iteration probes.  Grouped carry
rotation values expose mode-locking components, but we do not call their union
maximal.  A source-local factor \(1-z^n\lambda^n\) is recorded only as
itinerary bookkeeping.  There is no arithmetic clock or target determinant;
the strict Route-A verdict is \texttt{ROUTE\_A\_REJECTED}.
\end{abstract}

\section{Map, branches, and the question}
The fractional-part notation hides a genuine modelling choice at the jump.
We freeze the half-open state space \(I=[0,1)\) and write
\begin{equation}
 x_{j+1}=\lambda x_j+\delta-k_j,\qquad
 k_j\leq\lambda x_j+\delta<k_j+1,\qquad k_j\in\{0,1\}. \tag{1}
\end{equation}
Thus the lower carry equality belongs to a branch and the upper equality does
not.  The research question is whether this convention permits an exact
finite mode-locking certificate rather than a floating-point picture.  The
answer below is affirmative at an explicit cutoff, with all global claims
carefully withheld.

\section{Affine words and exact admissibility}
\begin{theorem}[Word return formula]
For \(w=(k_0,\ldots,k_{n-1})\),
\begin{equation}
 f_w(x)=\lambda^n x+\delta\sum_{r=0}^{n-1}\lambda^r
 -\sum_{j=0}^{n-1}k_j\lambda^{n-1-j}. \tag{2}
\end{equation}
Consequently the word has one candidate fixed point
\begin{equation}
 x_w(\delta)=\frac{\delta}{1-\lambda}
 -\frac{K_w}{1-\lambda^n},\qquad
 K_w=\sum_{j=0}^{n-1}k_j\lambda^{n-1-j}, \tag{3}
\end{equation}
and derivative \((f_w)'=\lambda^n\).
\end{theorem}
\begin{proof}
Substitution of (1) gives (2).  Since \(0<\lambda^n<1\), solving
\(f_w(x)=x\) gives (3), and differentiating (2) gives the derivative.
\end{proof}

Write \(x_j(\delta)=a_j\delta+b_j\), beginning with (3), and update
\((a_{j+1},b_{j+1})=(\lambda a_j+1,\lambda b_j-k_j)\).  The next result is
the computational core.
\begin{proposition}[Interval equivalence]
The parameter set for which \(w\) is an admissible cycle is exactly
\begin{equation}
 \mathcal D_w=\left\{\delta\in[0,1):
 0\leq a_j\delta+b_j<1,\quad
 k_j\leq\lambda(a_j\delta+b_j)+\delta<k_j+1\ \forall j\right\}. \tag{4}
\end{equation}
For rational \(\lambda\), this is an exactly computable rational interval,
possibly empty, with independently recorded endpoint closure.
\end{proposition}
\begin{proof}
The inequalities are necessary by the state domain and (1).  If they hold,
each selected carry is exactly \(k_j\), so (2) returns the candidate to itself.
Every inequality is an affine half-line in \(\delta\), and finite intersection
preserves exact rational endpoints and their open/closed status.
\end{proof}

\section{Primitive representatives and the finite atlas}
A word is primitive when it is not a repetition of a shorter block.  We keep
the lexicographically least cyclic rotation.  Its source-local rotation label
is
\begin{equation}
 \rho(w)=\frac{1}{n}\sum_{j=0}^{n-1}k_j. \tag{5}
\end{equation}
The exact factor \(1-z^n\lambda^n\) is useful for checking repetition and
derivative bookkeeping, but it is not a target determinant.

The census uses all primitive canonical words of lengths 1 through 12.  There
are 747 per slope and 2241 rows in total.  Only 138 rows have nonempty
\(\mathcal D_w\); grouped rows retain each component interval and its word id.
Table~\ref{tab:examples} gives representative rows (all omitted rows are in
the machine-readable receipt).
\begin{table}[t]
\centering
\caption{Exact word-certified components.  Every upper endpoint is open under
the convention (1).  These intervals are not asserted maximal.}
\label{tab:examples}
\begin{tabular}{ccccc}
\toprule
\(\lambda\) & word & \(\mathcal D_w\) & \(\rho(w)\) & \(\lambda^n\)\\
\midrule
\(1/2\) & 0 & \([0,1/2)\) & 0 & 1/2\\
\(1/2\) & 01 & \([2/3,5/6)\) & 1/2 & 1/4\\
\(1/2\) & 001 & \([4/7,9/14)\) & 1/3 & 1/8\\
\(2/3\) & 01 & \([3/5,11/15)\) & 1/2 & 4/9\\
\(3/4\) & 01 & \([4/7,19/28)\) & 1/2 & 9/16\\
\bottomrule
\end{tabular}
\end{table}

The endpoint audit evaluates every active equality, including the parameter
constraint \(\delta<1\).  For example, the lower endpoint \(2/3\) of the
\(\lambda=1/2\), word 01 component satisfies the carry-1 lower equality and is
accepted; at \(5/6\), carry-0 reaches its upper equality and is rejected.
This is why a decimal plot alone cannot certify the interval.

\section{Independent iteration and literature boundary}
For each slope we also probe eight base values \(\delta=i/8\) and every
distinct endpoint.  A separate implementation runs 360 iterations at 90
decimal digits, detects a repeated suffix of length at most 12, and compares
its affine fixed point with the exact rational candidate.  The 295-row ledger
is a control for the finite certificate, not a claim beyond the cutoff.

Laurent and Nogueira study this contracted-rotation family and its rotation
number, including algebraic-parameter phenomena~\cite{laurent2018}.  The
interval piecewise-contraction results of Nogueira and Pires give a finite
periodic-orbit bound under their stated hypotheses~\cite{nogueira2015}.  No
global one-periodic-orbit theorem is claimed here: their general two-branch
theorem gives only an at-most-two bound under its hypotheses.
Bugeaud and Conze supply the contracting-mod-one and Farey/Hecke--Mahler
context~\cite{bugeaud1999}.  Our contribution is the reproducible finite
Fraction census and boundary ledger, not a priority claim.

\ifnum\CRevisionRound>0
\section{Audit receipt and route boundary}
The independent checker validates all 2241 word rows and 295 direct rows,
SymPy validates 119 generic or rational identities, byte replay matches two
fresh producer runs, and 33 repaired-hash hostile mutations are rejected.
The three revision PDFs are content-distinct; each is built twice with the
fixed epoch \texttt{1788048000}, and the final PDF equals round 2 byte-for-byte.
\fi

\ifnum\CRevisionRound>1
The source has no intrinsic rational-prime carrier, so A0 fails.  A1 is limited
to the finite analytic word certificate.  A2 is an explicit target-match
failure: no target determinant or zero comparison is defined.  No continuation
or functional equation is supplied (A3\_FAIL), and the scalar contraction gives
only a formal lift hint (A4\_FORMAL\_HINT).  Thus
\[
(\mathtt{A0\_FAIL},\mathtt{A1\_PASS\_ANALYTIC},\mathtt{A2\_FAIL},
 \mathtt{A3\_FAIL},\mathtt{A4\_FORMAL\_HINT})
\]
and \texttt{ROUTE\_A\_REJECTED}; Route B is disabled.
\fi

\begin{thebibliography}{9}
\bibitem{laurent2018} M. Laurent and A. Nogueira, ``Rotation number of
contracted rotations,'' \emph{Journal of Modern Dynamics} 12 (2018),
175--191, DOI: \href{https://doi.org/10.3934/jmd.2018007}{10.3934/jmd.2018007}.
\bibitem{nogueira2015} A. Nogueira and B. Pires, ``Dynamics of piecewise
contractions of the interval,'' \emph{Ergodic Theory and Dynamical Systems}
35 (2015), 2198--2215, DOI:
\href{https://doi.org/10.1017/etds.2014.16}{10.1017/etds.2014.16}.
\bibitem{bugeaud1999} Y. Bugeaud and J.-P. Conze, ``Calcul de la dynamique
d'une classe de transformations linéaires contractantes mod 1 et arbre de
Farey,'' \emph{Acta Arithmetica} 88 (1999), 201--218, DOI:
\href{https://doi.org/10.4064/aa-88-3-201-218}{10.4064/aa-88-3-201-218}.
\end{thebibliography}

\section*{Declarations}
\textbf{Scope.} \texttt{NO\_BAD\_EULER\_OR\_ROOT\_NUMBER}; all determinant
language is source-local itinerary bookkeeping.  No target arithmetic,
target-zero, Euler-factor, root-data, automorphy, or Hilbert--Pólya claim is
made.  \textbf{Data and code.} Exact rows, checker, replay, and mutation audit
are included in the release.  This is not external peer review.
\textbf{AI-use disclosure.} Generative tools assisted drafting and code
generation; displayed formulas and metadata are checked by the deterministic
artifact chain.
\end{document}
